\begin{textsample}{POS Dim 1 – ac – Score 56.00 – ac/ac\_em\_18.txt}  \label{ex:f1_pos_171}
11.1.
Sistemáticas de acompanhamento e avaliação Na era da \textbf{tecnologia}, da globalização e dos avanços científicos em que \textbf{vivemos}, a Educação é \textbf{central} no processo de inserção e interação social dos sujeitos, portanto, no que \textbf{diz} respeito à sua qualidade de oferta, os processos \textbf{avaliativos} são de \textbf{suma} importante, pois \textbf{configura} -se como um dos pilares fundamentais do trabalho pedagógico da escola e da sala de \textbf{aula}.
Avaliação e aprendizagem devem caminhar juntas no processo educativo.
Dessa forma, não há avaliação sem aprendizagem, assim como não há aprendizagem sem avaliação.
Nessa mesma direção, Santos e Gontijo afirmam que : > [... ] avaliar é um processo natural que faz parte do repertório das ações dos seres humanos e está indissociavelmente ligado à busca de todo tipo de ação que provoque mudanças, mudanças essas que indicam respostas concretas ao desejo por “ algo melhor ”.
O ato de avaliar tem como \textbf{foco} a construção dos melhores resultados possíveis.
O ser humano permanentemente avalia em busca de uma melhor qualidade de vida, da construção de si mesmo e do melhor modo de ser e viver ( SANTOS ; GONTIJO, 2018, p. 31 ).
Portanto, a avaliação não deve ser somente o momento da realização das \textbf{provas} e testes, mas um processo contínuo e que ocorre dia após dia, visando a identificação dos objetivos de ensino alcançados, percebendo as dificuldades do \textbf{aluno} em seus movimentos para a constituição de conhecimentos e o domínio de habilidades, previstas para a sua etapa escolar.
Nesse sentido, a forma \textbf{avaliativa} funciona como um elemento de integração e motivação para o processo de \textbf{ensino-aprendizagem}.
A avaliação apresenta -se em várias modalidades, entre as quais está a avaliação somativa ou classificatória que, conforme Haydt ( 2000 ), tem como função classificar os \textbf{alunos} ao \textbf{final} da unidade, semestre ou ano letivo, segundo níveis de aproveitamento apresentados.
O objetivo da avaliação somativa é classificar o \textbf{aluno} para determinar se ele será aprovado ou reprovado e está vinculado à noção de medir.
O sistema educacional, muitas vezes, tem se apoiado na avaliação classificatória com a pretensão de \textbf{verificar} aprendizagem ou \textbf{competências} através de medidas, de quantificações.
Este tipo de avaliação pressupõe que as pessoas aprendam do mesmo modo, nos mesmos momentos e tentam evidenciar \textbf{competências} isoladas.
Ou seja, algumas pessoas, que por diversas \textbf{razões} têm maiores condições de aprender, aprendem mais e melhor.
Outras, com características diferentes, que não respondem tão bem ao conjunto de \textbf{disciplinas}, aprendem cada vez menos e são muitas vezes excluídas do processo de escolarização.
Dentre as modalidades de avaliação, há também outras duas que são : a avaliação formativa e a \textbf{diagnóstica}.
A avaliação formativa se constitui em um processo \textbf{complexo}, cujo detalhamento se dará na relação direta entre professores e estudantes, ou seja, deve ser dimensionada e modulada para cada realidade escolar, não havendo uma fórmula a ser aplicada, mas o desenvolvimento de um processo que é parte da própria aprendizagem do estudante e do fazer pedagógico do professor.
A referida modalidade de avaliação é chamada formativa no sentido de que indica como os \textbf{alunos} estão se \textbf{mobilizando} em direção aos objetivos visados.
Assim, a avaliação formativa tem como função informar o \textbf{aluno} e o professor sobre os resultados que estão sendo alcançados durante o desenvolvimento das atividades, \textbf{subsidiando} -se em melhorar o ensino e a aprendizagem ; localizar, apontar, discriminar deficiências, insuficiências, no desenvolvimento do ensino-aprendizagem para eliminá -las ; proporcionar feedback de ação ( leitura, explicações, exercícios ) ( SANT’ANNA, 2001, p. 34 ).
A avaliação formativa adota vários tipos de \textbf{métodos} \textbf{avaliativos}, sempre de \textbf{maneira} engajada, oferecendo ao \textbf{aluno} a oportunidade de ser coautor da construção do próprio conhecimento e, portanto, utiliza algumas ferramentas, como as seguintes : - Autoavaliação - Prerrogativa para estimular as \textbf{competências} \textbf{socioemocionais}, motivando a empatia e a sinceridade, estimula o protagonismo dos \textbf{alunos} e a prática da autogestão. - Testes tradicionais - O \textbf{erro} é visto como parte do processo e não como falta grave. - Simulados - Preparam os estudantes para determinados \textbf{exames}, estimulando um estudo com mais perseverança. - \textbf{Seminários} - Promove autonomia, protagonismo, estimula a participação ativa. - Trabalhos em grupo - Os \textbf{alunos} aprendem a lidar com a opinião do outro e com a diversidade e incentiva a empatia.
Vale \textbf{lembrar} que esse tipo de avaliação visa formar e informar os \textbf{alunos} a respeito de seu desenvolvimento como cidadãos cientes de seu lugar no mundo.
Portanto, os \textbf{métodos} \textbf{avaliativos} que constituem uma avaliação formativa podem ser abertos ao \textbf{método} didático do professor, mas é \textbf{imprescindível} que a autoavaliação esteja presente sempre.
Portanto, não há melhor ou \textbf{pior} avaliação ; cada avaliação possui objetivos distintos.
O \textbf{ideal}, no entanto, é usá -las em diferentes circunstâncias, unindo o melhor de cada modelo \textbf{avaliativo}.
As avaliações somativas são essenciais para informar e situar os estudantes da escola como um todo.
Todavia, a avaliação formativa é utilizada para um processo contínuo e longo em que o \textbf{erro} não é nada mais que um fator \textbf{inerente} ao processo de \textbf{ensino-aprendizagem} dos estudantes.
Outro perfil \textbf{avaliativo} é a avaliação \textbf{diagnóstica}, constituída por uma sondagem, \textbf{projeção} e retrospecção da situação de desenvolvimento do \textbf{aluno}, dando -lhe elementos para \textbf{verificar} o que aprendeu e como aprendeu.
É uma etapa do processo educacional que tem por objetivo \textbf{verificar} em que medida os conhecimentos anteriores ocorreram e o que se faz necessário planejar para selecionar dificuldades encontradas.
A avaliação \textbf{diagnóstica} contribui para a \textbf{compreensão} das particularidades de aprendizagem de cada indivíduo e orienta sobre o \textbf{direcionamento} mais assertivo sobre o que deverá ser feito a partir do que foi diagnosticado ; possui ação preventiva, promove o conhecimento de aptidões, interesses e \textbf{capacidades} para direcionar planejamentos futuros, revelando causas e dificuldades de aprendizagens.
Assim, seus resultados servem para identificar e propor novos procedimentos de \textbf{ensino-aprendizagem}.
Conforme a realidade escolar do docente, este tem a liberdade para realizar as diferentes modalidades de avaliação relatadas, fazendo isso no momento do planejamento de suas \textbf{aulas}.
A avaliação formativa possui um papel \textbf{central} no que se refere à organização do currículo para o ensino \textbf{médio}, \textbf{embasado} na repartição do tempo escolar em duas partes indissociáveis - Formação Geral Básica e Itinerários Formativos, ambas \textbf{articuladas} pelo Projeto de Vida.
Nesse sentido, a avaliação formativa atua como instrumento \textbf{metodológico} essencial para que sejam \textbf{atingidos} os objetivos de aprendizagem propostos no currículo.
Isso visa manter os pressupostos que levaram à elaboração deste documento e adequá -los, naquilo que for necessário, à nova legislação vigente.
Contudo, além de conhecer diferentes tipos de avaliação, é necessário também levar em conta os \textbf{métodos} de ensino que precedem as avaliações.
Nesse sentido, as metodologias ativas — práticas pedagógicas pensadas em conformidade com o que prega a BNCC — vem para modificar um pouco o ensino tradicional e dar base para que os estudantes tenham autonomia para trilhar o \textbf{caminho} rumo ao aprendizado, de \textbf{maneira} ativa.
As metodologias ativas, junto às avaliações formativa e somativa, são realmente interessantes e devem ser pensadas e desenvolvidas juntas, e não de \textbf{maneira} dissociada.
Nos últimos anos, o ensino \textbf{médio} tem passado por \textbf{inúmeras} mudanças, tendo seus marcos legais definidos por documentos norteadores, dentre os quais cabe destacar a Lei 13.415/2017 que, em seu Art. 3º, acrescentou o art. 35 -A, ampliando o texto da LDB 9394/96, e ressaltando ainda que : > § 7º Os currículos do ensino \textbf{médio} deverão considerar a formação integral do \textbf{aluno}, de \textbf{maneira} a adotar um trabalho voltado para a construção de seu projeto de vida e para sua formação nos aspectos físicos, cognitivos e \textbf{socioemocionais}. > § 8º Os conteúdos, as metodologias e as formas de avaliação processual e formativa serão organizados nas redes de ensino por meio de atividades \textbf{teóricas} e práticas, \textbf{provas} orais e escritas, \textbf{seminários}, projetos e atividades on-line, de tal forma que ao \textbf{final} do ensino \textbf{médio} o educando demonstre : > I - domínio dos princípios científicos e tecnológicos que presidem a produção \textbf{moderna} ; > II - conhecimento das formas \textbf{contemporâneas} de linguagem ( BRASIL, 2017, p. 2-3 ).
O Currículo de Referência Único do Acre propõe às escolas públicas da Rede de Educação Básica do Estado do Acre um processo \textbf{avaliativo}, orientado, normatizado e instrumentado pela instrução normativa nº 01 de 28 de fevereiro de 2019, \textbf{publicada} na \textbf{edição} nº 12.504 do D.O.E. de 01 de março de 2019, página 65, que considera o disposto na Lei de Diretrizes e Bases da Educação Nacional nº 9394/96, no Parecer do Conselho Nacional de Educação CNE/CEB nº 7/2010, na Resolução CNE/CEB nº 4/2010, na Resolução CNE/CEB nº 7/2010, no Parecer nº 15/2001 do Conselho Estadual de Educação – CEE e nas demais normas vigentes e que ainda serão elaboradas : > Art. 1{\EmojiFont °} Instruir o processo de avaliação da aprendizagem dos \textbf{alunos} do Ensino Fundamental e Médio das Escolas da Rede Pública de Educação do Estado do Acre face à interpretação dos instrumentos legais regulamentadores da matéria e assim estruturá -lo : > I. o ano letivo, para efeitos de organização e critérios que medem resultados, está dividido em 4 ( quatro ) bimestres ; > II. ao \textbf{final} de cada bimestre o \textbf{aluno} precisará ter alcançado uma nota \textbf{resultante} da aplicação de diversos instrumentos da avaliação formativa da aprendizagem tais como : \textbf{provas} e testes, trabalhos individuais ou em grupos, pesquisa, leituras complementares, \textbf{diagnósticos}, \textbf{tarefas} para a casa, exercícios desenvolvidos em sala de \textbf{aula}, tendo o alcance da média de qualidade que é 7,0 ( sete ).
A educação escolar orienta -se, intencionalmente, por metas que objetivam acompanhar todo o processo de ensino e aprendizagem.
Tais intenções da ação educativa só adquirem sentido quando levadas em consideração a natureza social e a função socializadora da educação escolar, as quais têm como \textbf{razão} \textbf{primordial} a promoção do desenvolvimento humano.
Incluem -se nesse processo os procedimentos didáticos assumidos pelo professor e a avaliação como ferramenta fundamental na aprendizagem ( DARSIE, 1996 ).
Como já comentado, avaliação e aprendizagem devem caminhar juntas.
Sendo assim, o professor, conforme a sua realidade escolar, poderá escolher as estratégias \textbf{avaliativas}, levando em consideração as propostas de atividades/metodológicas utilizadas para \textbf{mobilizar} o desenvolvimento de \textbf{competências} e habilidades, isto é, o professor poderá fazer uso de recursos \textbf{avaliativos} diferenciados.
Segundo Santos e Varela ( 2007 ), ao avaliar, o professor deve utilizar técnicas diversas e instrumentos variados, para que assim consiga diagnosticar o começo, o durante e o fim de todo o processo \textbf{avaliativo} percorrido pelo \textbf{discente}, para que, a partir de então, possa progredir no processo didático e retomar o que foi insatisfatório para o processo de aprendizagem dos educandos.
Por fim, os estudantes de ensino \textbf{médio}, para alcançarem determinado aprendizado, deparam -se com um processo \textbf{complexo}, de cunho pessoal, para o qual os professores e a escola podem contribuir ao permitir a comunicabilidade, ao situá -los em seu grupo e a oportunizar que aprendam a respeitar e a fazer -se respeitar.
Para tal, é salutar que professores e escolas criem situações em que o estudante deva sentir -se desafiado ou \textbf{instigado} a participar e questionado pelo jogo do conhecimento, adquirindo o espírito de pesquisa e o desenvolvimento de \textbf{capacidades} de \textbf{raciocínio} e de autonomia ( BRASIL, 1999 ).
Assim, os \textbf{parâmetros} de avaliação voltados para a etapa do ensino \textbf{médio} devem estar em consonância com os princípios que \textbf{permeiam} esta etapa de ensino, de modo a \textbf{subsidiar} o professor para diagnosticar e monitorar as aprendizagens, identificar as potencialidades e eventuais dificuldades presentes no processo, além de acompanhar os resultados das práticas de ensino.
Desse modo, avaliação deverá ter como referência : - O desenvolvimento integral do \textbf{aluno}, de modo a garantir a autonomia e o apropriação das \textbf{competências} \textbf{socioemocionais}. - Flexibilização de metodologias para efetivo desenvolvimento das \textbf{competências} cognitivas e \textbf{socioemocionais}. - Coerência entre a prática pedagógica e os processos \textbf{avaliativos}. - A avaliação dos resultados das avaliações da aprendizagem, de modo a proporcionar reflexões acerca das evidências e \textbf{diagnósticos} para os ajustes necessários nas práticas pedagógicas. - A avaliação formativa deve está constantemente presente como instrumento, de modo a garantir a aferição do desenvolvimento das \textbf{competências} gerais e específicas do Currículo. - Os instrumentos \textbf{avaliativos} devem ser diversificados e articulados entre si. - As avaliações devem acompanhar a flexibilização do ensino \textbf{médio} e apoiar os jovens em seus Projetos de Vida.

% matched lemmas: aluno, articulado, atingir, aula, avaliativo, caminho, capacidade, central, competência, complexo, compreensão, configurar, contemporâneo, diagnóstico, direcionamento, discente, disciplina, dizer, edição, embasar, ensino-aprendizagem, erro, exame, final, foco, ideal, imprescindível, inerente, instigar, inúmero, lembrar, maneira, mau, metodológico, mobilizar, moderno, médio, método, parâmetro, permear, primordial, projeção, prova, publicar, raciocínio, razão, resultante, seminário, socioemocional, subsidiar, suma, tarefa, tecnologia, teórico, verificar, vivar
\end{textsample}
